\section{Methods under study}\label{sec:methods}

\begin{figure}[t]
\centering
\figorplaceholder[\linewidth]{figs/f2_method.pdf}{2.37in}
\caption{\textbf{The RoPE base changes frequencies, a global temperature scales all terms, and per-band gains reweight individual terms.} (a)~Query, key and offset on the token grid of a $4096^2$ image; the inset shows band 8 under the native base $b$ (dashed) and \sega's released base $b'$. (b)~Four of the \N{c@pairsperaxis} width-axis bands with their total log gains (sky: $\ln\tau$; circles: \segaofficial; diamonds: \constancytm) from one held-out sample at step 14. The three separate traces sum these four bands with unit content amplitudes and zero phases on the same vertical scale; they illustrate the base and temperature changes, not measured query--key logits. (c)~Two alternative controllers set the total gain $g_j$, which already includes temperature, from their respective spectral measurements.}
\label{fig:method}
\end{figure}

\begin{table}[t]
\centering
\caption{\textbf{Arms and the knobs they turn.} Temperature is the logit gain at $4096^2$; t/x: varies with the noise level / with the image. Held-out arms ran on rows 8--194, the pilot and dev arms only on the eight dev prompts.}
\label{tab:methods}
\footnotesize
\setlength{\tabcolsep}{2.6pt}
\begin{tabular}{@{}l l c c c l@{}}
\toprule
Arm & Frequency base & Temperature & Per-band gain & t/x & $n$ \\
\midrule
\ntk & $\theta\,s^{2D/(D-2)}$ & none & none & no / no & \N{ho@nprompts} \\
\dype & PI/NTK ramps in $\sigma$ & $\N{c@dypegain}$ (spatial) & none & yes / no & \N{ho@nprompts} \\
\segaofficial & $\theta\,s^{2D/(D-2)}/(1{+}0.1\ln s)$ & $s^{0.16}=\N{c@segatemp}$ & sensor, $\ge1$ & yes / yes & \N{ho@nprompts} \\
\segapaper & $\theta\,s^{2D/(D-2)}$ & $s^{0.16}$ & sensor, unbounded & yes / yes & 8 (pilot) \\
\tempsb & as \segaofficial & $s^{0.16}$ on all dims & none & no / no & \N{ho@nprompts} \\
\constancy & as \ntk & $\tau_N=\N{c@taulogn}$ & gated, late target & yes / yes & \hon{cons} (partial) \\
\constancytm & as \segaofficial & $s^{0.16}$ & gated, time-matched & yes / yes & \N{ho@nprompts} \\
\aliastemp & as \segaofficial & none & alias count, static & no / no & 8 (dev) \\
\segactm & as \segaofficial & $s^{0.16}$ & \sega{} $\times$ time-matched & yes / yes & \N{ho@nprompts} \\
\bottomrule
\end{tabular}
\end{table}



All arms share the sampler and differ only in $\omega'$ and $g$ of \cref{eq:fourier} (\cref{tab:methods}; exact formulas and code revisions in \cref{app:impl}).

\paragraph{Frequency baselines}
\ntk{} is the strengthened NTK base that \sega{} builds on, $\theta'=\theta s^{2D/(D-2)}$, which divides the slowest pair's frequency by $s^2$. With no temperature and no allocation, it is the floor of the decomposition. \dype~\citep{issachar2025dype}, ported from the official code, blends position interpolation, NTK and native frequencies with two ramps in $\sigma^2$, so fine pairs return to their native frequencies late in denoising; its cos/sin amplitude $1+0.1\ln s$ acts as a logit gain of $\N{c@dypegain}$ on the spatial dimensions.

\paragraph{\sega: paper equations and released code}
\sega~\citep{sega2026} sets per-pair gains from its input spectrum around $m_{\mathrm{ref}}=s^{0.08}$, assigning more weight to low-energy pairs. Our reading of the printed formulation differs from the release at each of six compared stages (\cref{tab:papercode,app:impl}); we call the parity-tested released port \segaofficial{} and the printed-law implementation \segapaper.

\paragraph{\tempsb: the decomposition control}
\tempsb{} keeps the frequency base of \segaofficial{} and applies its reference logit gain $s^{0.16}$ uniformly to every dimension. It differs from \segaofficial{} in three ways: no per-band shape, no lift in mean gain from the clamp at 1, and a gain of $s^{0.16}$ rather than 1 on the index axis.

\subsection{Per-band laws from vision science}\label{sec:ours}

Contrast constancy describes nearly stable perceived band contrast above threshold across viewing distances~\citep{georgeson1975contrast}. Efficient coding combines visibility against noise, $S/(S+N)$, with contrast normalization toward natural-image statistics~\citep{atick1990towards,atick1992retina,vanhateren1992theory,field1987relations,webster2002neural,parraga2000human}. We test four controllers inspired by these ideas; full equations and calibration details are in \cref{app:laws}.
\begin{itemize}
\item \textbf{\constancy.} A noise-visibility gate weights a bounded log-power correction toward a native-resolution target from the final clean steps. The law uses the \ntk{} base and log-$N$ temperature.
\item \textbf{\constancytm.} The same correction uses a target matched to the current noise level and the base and temperature of \segaofficial{}, with \tempsb{} as its control.
\item \textbf{\aliastemp.} A pair's resolvable key count sets its log-$N$ temperature; \aliastemps{} shifts the log gains to match \sega's reference geometric mean.
\item \textbf{\segactm.} The time-matched correction multiplies \sega's released allocation, $\ln g=\ln g^{\textsc{sega}}+\Delta$. Development selection chose $\eta=1$ for the held-out run.
\end{itemize}
